% =====================================================================
% Appendix E -- Intended versus Executed
% =====================================================================
\documentclass[../preamble.tex]{subfiles}
\begin{document}

\chapter{Intended versus Executed}
\label{app:intended-executed}

\section{Motivation}
\label{app:intended-why}

The pipeline design and the executed experiment were not identical. Several
features were specified in the v3 methodology document but not used in the
final reported run. This appendix makes the differences explicit so the
thesis can be read as a record of what was actually done, not as a record
of what was intended.

\section{Differences table}
\label{app:intended-table}

\begin{table}[htbp]
  \centering
  \small
  \caption{Intended design versus executed run. The ``Source of truth'' column
           names the artefact that justifies the executed value.}
  \label{tab:intended-vs-executed}
  \begin{tabularx}{\linewidth}{@{}lXXX@{}}
    \toprule
    \textbf{Item} & \textbf{Intended (v3 spec)} & \textbf{Executed} & \textbf{Source of truth} \\
    \midrule
    Default freeze fraction       & $0.60$, HPO in $\{0.40,0.50,0.60,0.70\}$
                                   & $0$, HPO in $\{0, 0.5\}$
                                   & HPO best-params YAML files \\
    EMA                           & $\beta = 0.999$, enabled
                                   & not enabled
                                   & training logs (no EMA shadow state) \\
    SWA                           & over last 5 checkpoints
                                   & not used
                                   & no SWA buffer in checkpoint files \\
    $k$-fold cross-validation     & $k = 5$ over combined corpus
                                   & single 80/10/10 split
                                   & sweep matrix CSV \\
    Search-space dropout          & $[0.3, 0.7]$
                                   & $[0, 0.5]$
                                   & \texttt{experiment\_config.yaml} \\
    Search-space LR               & $[1 \times 10^{-5}, 1 \times 10^{-3}]$
                                   & $[5 \times 10^{-6}, 1 \times 10^{-4}]$
                                   & \texttt{experiment\_config.yaml} \\
    Search-space WD               & $[1 \times 10^{-4}, 1 \times 10^{-1}]$
                                   & $[1 \times 10^{-5}, 1 \times 10^{-2}]$
                                   & \texttt{experiment\_config.yaml} \\
    Rotation augmentation        & not in spec
                                   & $[0, 20]$ degrees
                                   & \texttt{experiment\_config.yaml} \\
    Label smoothing              & not in spec
                                   & $[0, 0.1]$
                                   & \texttt{experiment\_config.yaml} \\
    XAI                           & Grad-CAM, LRP, SHAP
                                   & Grad-CAM only
                                   & thesis scope evolution \\
    Calibration                   & single-pass temperature scaling
                                   & two-pass temperature scaling
                                   & calibration JSON \\
    External cohort              & not in spec
                                   & PCOSgen (Sundari et~al., 2025)
                                   & dataset inventory \\
    \bottomrule
  \end{tabularx}
\end{table}

\section{Why the deviations exist}
\label{app:intended-why-deviations}

\subsection*{Backbone freezing}

The v3 spec assumed that freezing the lower 60\% of the backbone would
prevent overfitting on the small source dataset. The empirical evidence was
the opposite: when the search space was widened to include full
fine-tuning, Optuna consistently selected \texttt{freeze\_fraction=0} for
the architectures that contributed to the final ensemble. This is consistent
with the observation that the external cohort introduces new feature
statistics that benefit from full adaptation.

\subsection*{EMA, SWA, and $k$-fold}

These features were implemented in the trainer scaffold but were not used
to produce the reported results. EMA and SWA were deprioritised because the
calibration step already improved the calibration error and the additional
engineering effort was not justified by the expected gain on the small
validation set. $k$-fold cross-validation was deprioritised because
PCOSgen does not document a verified patient-level grouping, so a true
$k$-fold scheme would require a unified dataset with new metadata.

\subsection*{Search-space bounds}

The v3 spec gave wider dropout, learning-rate, and weight-decay bounds.
The executed search space is narrower, which makes the HPO step more
tractable on a single GPU and produces reproducible best parameters. The
chosen bounds were validated by a small preliminary sweep that confirmed
the optimum lay within the narrower ranges.

\subsection*{Rotation and label smoothing}

Rotation and label smoothing were not in the v3 spec but were added in the
executed search space because they appear commonly in medical-imaging HPO
studies. They were retained because at least one ensemble member selected
a non-zero value (rotation=13 for ConvNeXt-Tiny, label smoothing 0.092 for
ViT-B/16, label smoothing 0.030 for ConvNeXt-Tiny).

\subsection*{Two-pass calibration}

The v3 spec described a single-pass temperature scaling. The executed
pipeline applies per-model temperature first, then a single ensemble
temperature on the averaged probabilities. The two-pass approach was
adopted after observing that the per-model and ensemble temperatures
differed substantially (e.g., ensemble $T = 0.93$ versus per-model
$T \in \{2.02, 1.96, 2.20\}$).

\subsection*{External cohort}

PCOSgen was added after the v3 spec was written because it became publicly
available. The thesis treats PCOSgen as an external validation cohort, not
as additional training data, so the addition strengthens rather than
redesigns the experiment.

\end{document}